\subsubsection{Design of a RL-Based HAC}

Alternatively, reinforcement learning (RL) can also be used in building an HAC agent with Unity's Machine Learning Agents (ML-Agents). ML-Agents allows developers to utilize games and simulations as training environments for intelligent agents. These agents can be trained using various paradigms, including reinforcement learning, imitation learning, and neuroevolution. Practical applications of trained agents encompass the control of non-player character (NPC) behavior in multi-agent or adversarial settings, automated game build testing, and the pre-release evaluation of game design decisions.

To facilitate training, a configuration file is typically required for the Python API, with Proximal Policy Optimization (PPO) serving as the predominant algorithm. The Python API interacts with the Unity environment via a communicator to execute the training process. Upon completion, the trained model is exported as an Open Neural Network Exchange (ONNX) file, which is subsequently used for agent inference and control. Furthermore, Unity provides a heuristic API for debugging and imitation learning (IL), enabling manual control over agents and the adjustment of training outcomes, as illustrated in \cref{fig:agent}.

\begin{figure}
    \caption{Simplified ML-Agents Scene Block Diagram}\label{fig:agent}
    \centering
\begin{tikzpicture}[
    font=\sffamily\bfseries,
    base/.style={rectangle, rounded corners=5pt, minimum width=2.2cm, minimum height=1cm, align=center, draw=white, line width=1pt, text=white},
    agent/.style={base, fill={rgb,255:red,119; green,184; blue,103}},
    behavior/.style={base, fill={rgb,255:red,98; green,155; blue,210}, minimum width=3cm},
    purple_box/.style={base, fill={rgb,255:red,146; green,106; blue,196}},
    orange_box/.style={base, fill={rgb,255:red,243; green,156; blue,87}},
    red_box/.style={base, fill={rgb,255:red,219; green,97; blue,94}},
    yellow_box/.style={base, fill={rgb,255:red,241; green,190; blue,47}},
    arrow/.style={<->, >=Stealth, thick, white},
    % 使用 * 并在连线时稍微偏移
    dash_arrow/.style={*-, dashed, thick, black} 
]

    % --- 内部节点 ---
    \node[behavior] (behA) {Behavior A};
    \node[agent, above left=1.5cm and -1.2cm of behA] (a1) {Agent\\A1};
    \node[agent, above right=1.5cm and -1.2cm of behA] (a2) {Agent\\A2};
    \node[orange_box, below=1cm of behA] (comm) {Communicator};

    \node[behavior, right=1cm of behA, minimum width=2.2cm] (behC) {Behavior C};
    \node[agent, at={(behC |- a1)}] (c1) {Agent\\C1};
    \node[purple_box, below=1cm of behC] (inf) {Inference};

    \node[behavior, right=1cm of behC, minimum width=2.2cm] (behD) {Behavior D};
    \node[agent, at={(behD |- a1)}] (d1) {Agent\\D1};
    \node[purple_box, below=1cm of behD] (heu) {Heuristic};

    % --- 外部节点 ---
    \node[red_box, below=2cm of comm] (pyAPI) {Python API};
    \node[yellow_box, right=1.2cm of pyAPI] (pyTrain) {Python Trainer};

    % --- 连线 ---
    \draw[arrow] (behA.north) -- ++(0,0.6) -| (a1.south);
    \draw[arrow] (behA.north) -- ++(0,0.6) -| (a2.south);
    \draw[arrow] (behA) -- (comm);
    \draw[arrow] (c1) -- (behC);
    \draw[arrow] (behC) -- (inf);
    \draw[arrow] (d1) -- (behD);
    \draw[arrow] (behD) -- (heu);
    
    % 虚线连线：起点从 comm.south 稍微往下移 1pt，确保圆点（*）不被框线切掉
    \draw[dash_arrow] ([yshift=-1pt]comm.south) -- (pyAPI.north);
    \draw[arrow, black] (pyAPI) -- (pyTrain);

    % --- 背景框 ---
    \begin{scope}[on background layer]
        % 增加 inner ysep=30pt 确保上下都有足够的留白
        \node[fill=gray, rounded corners=10pt, 
              fit=(a1) (a2) (c1) (d1) (comm) (inf) (heu), 
              inner xsep=15pt, inner ysep=30pt] (env) {};
        
        \node[anchor=north west, text=white, inner sep=10pt] at (env.north west) {Learning Environment};
    \end{scope}

\end{tikzpicture}
\end{figure}


For this research, each Drone Agent includes, (1) Agent Script, Handles OnActionReceived (processing joystick inputs), CollectObservations (monitoring states), and Heuristic (manual control for debugging); (2) Behavior Parameters: Defines the vector space for actions and observations. (3) Decision Requester: Controls the frequency of actions. By limiting the request rate, we can simulate more human-like reaction times, making the agent a better tutor for human pilots.

Proximal Policy Optimization (PPO) is an on-policy actor-critic algorithm that updates policies in small, clipped steps to ensure stability \parencite{openai2017ppo}. PPO is particularly effective for drones learning to navigate complex environments \parencite{yu2025depth}.

The optimization objective is defined as \cref{eq:ppo} shows.

\begin{equation}\label{eq:ppo}
L^{CLIP}(\theta) = \hat{\mathbb{E}}_t [ \min(r_t(\theta)\hat{A}_t, \text{clip}(r_t(\theta), 1-\epsilon, 1+\epsilon)\hat{A}_t) ]
\end{equation}

\begin{itemize}
    \item $\theta$ is the policy parameter
    \item $\hat{E}_t$ denotes the empirical expectation over timesteps
    \item $r_t$ is the ratio of the probability under the new and old policies, respectively
    \item $\hat{A}_t$ is the estimated advantage at time $t$
    \item $\epsilon$ is a hyperparameter, usually $0.1$ or $0.2$
\end{itemize}

Compared with a PID AI agent, a reinforcement learning-based AI agent does not require repeated testing to fine-tune parameters. Decisions are executed through a Decision Requester, which makes it easier to adjust the request frequency and pacing. However, as a black-box approach, it can be more difficult to achieve optimal performance or interpret and refine its behavior.